% #############################################################################
% This is Chapter 1
% !TEX root = ../main.tex
% #############################################################################
% Change the Name of the Chapter i the following line
\fancychapter{Introduction}
% The following line allows to ref this chapter
\label{chap:intro}
\noindent 
Language modeling seeks to estimate the likelihood of word sequences and to use context to predict or generate text. 
When such models are trained at very large scale, they enter the regime of Large Language Models (LLMs), which have become strong general-purpose tools
for tasks such as classification, generation, translation, and structured text
transformation \citep{zhao2025surveylargelanguagemodels}. Their success, however, should not hide the
basic limitation their behavior being shaped by the
distribution of the data on which it was trained. For languages with multiple
standardized varieties, models trained on generic web-scale data often become
better aligned with the dominant variety and less reliable for the others.

This issue is not specific to Portuguese. Many widely spoken languages are
pluricentric, in the sense that they are shared across different countries and
speech communities. Geographic separation, migration, institutional standardization, media
ecosystems, and local cultural practices all contribute to differences in
spelling, morphosyntax, lexical choice, and discourse conventions.
Comparable challenges arise in languages such as English, Spanish, Arabic, or
French, where users may expect computational systems to respect the variety
they actually use rather than the variety most represented in the training data. From an NLP perspective, these differences are
large enough to affect downstream performance, yet subtle enough that they are
frequently ignored during data collection and pre-training.

The Portuguese case makes this problem especially visible. Portuguese is the
official language of eight countries and is spoken by approximately 200 million of
people worldwide. Within this linguistic
space, European Portuguese (EP) and Brazilian Portuguese (BP) are the two most
prominent standardized varieties. While some differences are orthographic, many are
lexical, and others involve morphosyntactic preferences such as clitic
placement, verbal periphrases, subject realization, or forms of address \citep{barreiro-mota-2018-paraphrastic}.

The challenge is not only linguistic, but also sociotechnical. Brazilian Portuguese
overwhelmingly dominates the amount of Portuguese text available on the web, and
resources labeled simply as Portuguese tend to overrepresent Brazilian
usage. That imbalance propagates to pre-trained models, which are more likely to treat BP as the unmarked default
and to under-serve EP users in generation, classification, and translation
settings \citep{sousa2025enhancingportuguesevarietyidentification}. The problem can become even more severe when automatic
filtering pipelines used to build large training corpora discard material that
looks noisy, unusual, or weakly aligned, since dialectal and sociolectal text
is often disproportionately affected by such heuristics
\citep{sousa2025enhancingportuguesevarietyidentification}. At the same
time, although commercial systems increasingly expose language-variant options,
direct support for EP-BP rewriting remains much more limited than translation
between clearly different languages.

Research on NLP for language varieties has grown in recent years, but it still
remains comparatively limited when the goal is to jointly address variety
identification, translation between varieties, user-controlled generation, and
the data-quality issues that shape all of these tasks. This work aims to take
Portuguese as a case study for that broader problem. Rather than
treating Portuguese as a single homogeneous language, it explicitly models
intra-language variation and studies how corpus construction, data filtering,
model architecture, and fine-tuning strategies influence variety-aware behavior.

% The work reported in this dissertation developed along two closely connected
% axes. The first is data-centric: consolidating scattered EP-BP resources into a
% single corpus and subjecting them to a common preprocessing and quality-control
% pipeline. The second is model-centric: training and comparing systems for both
% variety identification and rewriting, with the goal of understanding which
% modeling choices matter most under realistic data and compute constraints.

This research project makes four main contributions:

\begin{enumerate}
    \item It consolidates a reusable EP-BP corpus by combining resources such
    as OPUS OpenSubtitles \citep{4992de1b5fb34f3e9691772606b36edf}, FRMT \citep{riley2023frmtbenchmarkfewshotregionaware}, the Gold Collection \citep{sanches2024brazilianportugueseeuropeanportuguese}, and
    PtBrVarId \citep{sousa2025enhancingportuguesevarietyidentification}, together with a documented preprocessing and filtering
    pipeline tailored to the quality issues of each source.
    \item It develops and evaluates models for both language variety
    identification and bidirectional rewriting between EP and BP.
    \item It provides a comparative study of model architecture and adaptation
    strategy, with particular emphasis on encoder-decoder models at two scales and on the trade-off between full fine-tuning and
    parameter-efficient adaptation.
    \item It complements aggregate evaluation with corpus analysis and
    qualitative error analysis, highlighting where the data is strong, where
    the models fail, and which linguistic phenomena remain difficult even after
    fine-tuning.
\end{enumerate}

The corpus contribution goes beyond simply concatenating existing datasets. The
resulting resource combines parallel data, variety-labeled data, and
automatically derived pairs, and it is accompanied by a preprocessing pipeline
that includes normalization, sentence merging when needed, boilerplate removal,
deduplication, length-based filtering, and alignment-based quality control. This is especially important for subtitle-derived data, where
segmentation noise, misalignment, and domain-specific artifacts can otherwise
dominate the signal learned by the model. Treating corpus construction as a
central research component, rather than as a preliminary cleanup step, is
therefore part of the contribution of this thesis.

From the modeling side, we study encoder-decoder adaptation
for classifying whether a sentence is EP or BP, and rewriting text from one variety
into the other while preserving meaning. The main model usede is
T5Gemma2 \citep{zhang2025t5gemma2seeingreading}, explored at 270 million and 4 billion parameters.
Given the available compute budget and the practical relevance of efficient
adaptation, the larger configuration is trained primarily with Low-Rank
Adaptation (LoRA) adapters \citep{hu2021loralowrankadaptationlarge}, while the smaller configuration is
fully fine-tuned to provide a strong comparison point \citep{houlsby2019parameterefficienttransferlearningnlp}. The results
obtained throughout this work show that carefully curated data and targeted
fine-tuning can yield systems that slightly outperform previously tested baselines for
EP/BP classification and rewriting.


The remainder of this dissertation is organized as follows. \textbf{FINISH LATER}
